% CHECKPOINT36: frozen elliptic_boundary_power_automorphy35.md
% SHA256 aa4a7809b09028d8e09183ee257dcbadbd535d011833f3e9d8fcd107e15c7c74.
% Includes the separately completely reviewed free-factor restriction corollary.
\statement{A.123}{Proved}{A boundary-power automorphy criterion and the original free-factor restriction}
Let \(E=\widehat F(X,Y)\), \(K=[Y,X]\), and let \(\eta\) be a
continuous endomorphism. If
\[
 \eta(K)=K^d,\qquad d\in\widehat{\mathbb Z}^{\times},            \tag{A.123.1}
\]
then \(\eta\) is an automorphism of \(E\).
For the actual candidate \(\eta\) of A.122, its unproved boundary
equation would therefore imply both
\[
 \Theta=\eta c_{\mathrm{ell}}\eta^{-1}
 \quad\text{on }E,\qquad
 \psi_d\in\operatorname{Aut}\bigl(\widehat F(p,q)\bigr).         \tag{A.123.2}
\]
The latter conclusion is about the actual original free factor
and its original normalized \(\psi_d\).

\begin{proof}
\emph{1. An actual finite-cover homology fact.}
Identify \(F_0=F(X,Y)\) with the fundamental group of a compact
orientable once-punctured torus \(S\), choosing its orientation
and a boundary whisker so that its boundary word is \(K\).
Let \(V\) be a proper open subgroup of \(E\) containing \(K\),
and put \(V_0=V\cap F_0\).
Density on the finite coset set gives
\([F_0:V_0]=[E:V]=n>1\).
The connected \(n\)-sheeted cover \(S_0\) is a compact orientable
surface with boundary.

Its boundary components are the orbits of \(K\) on the finite
coset set, with covering degree equal to orbit size.
Since \(K\in V_0\), the base coset is a singleton orbit.
It gives a degree-one boundary component whose based word is
the actual \(K\). Since \(n>1\), another orbit remains, and
the number \(b\) of boundary components is at least two.

Cutting a compact orientable genus-\(g\) surface with \(b\) boundary
components along handle arcs and arcs joining its boundaries
gives the polygon presentation
\[
 \left\langle
 a_1,b_1,\ldots,a_g,b_g,c_1,\ldots,c_b
 \ \middle|\
 \prod_{j=1}^g[a_j,b_j]c_1\cdots c_b=1
 \right\rangle .
\]
Its integral abelianization has the independent handle classes
and the boundary classes with the sole relation
\(c_1+\cdots+c_b=0\). For \(b\geq2\), each boundary class is
nonzero and primitive. Thus the class of \(K\) in
\((V_0)_{\mathrm{ab}}\) is nonzero and primitive.

This concerns the actual topology. Every finite-index subgroup
of \(V_0\) has finite index in \(F_0\), and its core in \(F_0\)
is finite-index and normal. Therefore the full profinite topology
of \(V_0\) is exactly the induced topology, and
\(\widehat V_0=V\). The universal finite-abelian quotient
property, or a free basis of \(V_0\), then gives
\[
 V_{\mathrm{ab}}
   =\widehat{(V_0)_{\mathrm{ab}}}
   \simeq\widehat{\mathbb Z}^{\,2g+b-1}.
 \]
The integral primitive class of \(K\) remains nonzero there.
No completion theorem for arbitrary closed surface subgroups
is used.

\emph{2. Surjectivity of the actual endomorphism.}
The image \(M=\eta(E)\) is compact and hence closed.
Condition (A.123.1) and the unit property give
\(K=(K^d)^{d^{-1}}\in M\).
If \(M\ne E\), choose \(t\notin M\). Compactness and finite
separation give an open normal \(U\) with \(t\notin MU\).
Then \(V=MU\) is a proper open subgroup containing \(M\)
and \(K\), so the preceding finite-cover argument applies.

Both \(\eta(X)\) and \(\eta(Y)\) lie in \(V\), whereas
\[
 K^d=[\eta(Y),\eta(X)]\in\overline{[V,V]}.
\]
Thus \(d[K]=0\) in \(V_{\mathrm{ab}}\). Multiplication by
the unit \(d\) is invertible on this free profinite abelian
group, contradicting the nonzero class just proved.
Consequently \(\eta\) is onto.

\emph{3. The Hopfian step.}
A finitely generated profinite group has finitely many open
normal subgroups of any fixed index. For a surjective continuous
endomorphism, inverse image is an injection on that finite set
and preserves index, hence is a permutation.
For each open normal \(M_0\) there is therefore an open normal
\(M_1\) with \(\eta^{-1}(M_1)=M_0\).
The kernel of \(\eta\) is contained in every such \(M_0\);
it is trivial. Hence \(\eta\) is an automorphism.

\emph{4. Restriction to the actual original free factor.}
Retain the candidate of A.122, and assume its exact boundary
equation. We have just proved \(\eta\in\operatorname{Aut}(E)\).
Since \(U,V\) are commutator-word evaluations and \(p=[Y,X]\),
its abelianization is \(\operatorname{diag}(d,1)\).
In particular the \(Y\)-parity character is unchanged, so
\(\eta(V_2)=V_2\), where
\[
 V_2=\ker(E\longrightarrow C_2),\qquad
 H=\overline{\langle p,q\rangle}
\]
are the actual groups of A.122.
That proof gives \(V_2=H\amalg\widehat{\mathbb Z}\) in its
explicit Schreier marking. Thus \(\eta(H)\) is also a free
factor of \(V_2\).

The assumed boundary equation gives \(\eta(p)=p^d\), while
\(\eta(q)=u(p,q)^{-1}q^du(p,q)\).
Consequently \(\eta(H)=\psi_d(H)\subseteq H\).
Let \(r:V_2\to\eta(H)\) be the free-factor retraction.
Its restriction \(r|_H\) is onto because it fixes
\(\eta(H)\subseteq H\). The continuous map
\[
 (\eta|_H)^{-1}\circ r|_H:H\longrightarrow H
\]
is therefore a surjective endomorphism of the finitely generated
profinite group \(H\). By the preceding Hopfian argument it is
injective, so \(r|_H\) is injective. An injective retraction
onto a subgroup has equal domain and image: for \(h\in H\),
\(r(h)=r(r(h))\), and injectivity gives \(h=r(h)\).
Hence \(\eta(H)=H\). Its restriction is exactly the original
\(\psi_d\), proving the second assertion in (A.123.2).
The first follows from the already proved A.122 equivalence
\(\Theta\eta=\eta c_{\mathrm{ell}}\).

The boundary equation remains a hypothesis. This proof derives
neither that equation nor the original full hexagon or CM sign.
It uses actual finite covers and the displayed free factor,
rather than a transfer of pro-\(p\) automorphy or an unverified
discrete-retract assertion.
\end{proof}
